% !TeX spellcheck = fr_FR
\PassOptionsToPackage{hidelinks}{hyperref}

\documentclass[10pt,a4paper,openany]{book} % openany: Kapitel starten auf beliebiger Seite

% --- Sprache & Encoding ---
\usepackage[T1]{fontenc}
\usepackage[utf8]{inputenc}
\usepackage[english]{babel}

% === SCHRIFTEN (sichtbar anderes Design) ======================
\usepackage{tgtermes}        % Serif (Text) Times-ähnlich
\usepackage{newtxmath}       % Mathe passend zu Times
\usepackage[scaled=0.92]{helvet} % Sans-Serif (Überschriften)
\usepackage{inconsolata}     % Monospace (Code)
% ===============================================================

% --- Seitenränder kompakt ---
\usepackage[a4paper,margin=1.8cm]{geometry}

% --- Grundpakete ---
\usepackage{graphicx}
\usepackage{svg}
\svgsetup{inkscapelatex=true}
\usepackage{pdfpages}
\usepackage{booktabs}
\usepackage{array}
\usepackage{multirow}
\usepackage{ltablex}\keepXColumns
\usepackage{siunitx}
\usepackage[printonlyused,withpage]{acronym}
\usepackage[numbers]{natbib}
\usepackage{fancyhdr}
\usepackage{comment}
\usepackage{rotating}
\usepackage{import}
\usepackage{pythontex}
\usepackage{adjustbox}
\usepackage{float}
\usepackage{tikz}
\usepackage{hyperref}

% --- Listings: farbiger Code ---
\usepackage{xcolor}
\usepackage{listings}

\definecolor{codegreen}{rgb}{0,0.6,0}
\definecolor{codegray}{rgb}{0.5,0.5,0.5}
\definecolor{codepurple}{rgb}{0.58,0,0.82}
\definecolor{backcolour}{rgb}{0.95,0.95,0.92}

\lstdefinestyle{mystyle}{
  backgroundcolor=\color{backcolour},
  commentstyle=\color{codegreen},
  keywordstyle=\color{magenta},
  numberstyle=\tiny\color{codegray},
  stringstyle=\color{codepurple},
  basicstyle=\ttfamily\footnotesize,  % nutzt Inconsolata
  breaklines=true,
  numbers=left,
  numbersep=5pt,
  showstringspaces=false,
  tabsize=2,
  frame=single,
  captionpos=b
}
\lstset{style=mystyle}

% --- Überschriften in Sans-Serif (Helvetica, fett) ---
\usepackage{sectsty}
\allsectionsfont{\sffamily\bfseries}

% ===== keine automatisch eingefügten Leerseiten =====
\let\cleardoublepage\clearpage
% ====================================================

% --- Deine eigenen Kommandos & Config ---
%
% own commands
%

%double empty page
\newcommand \myemptypage {
    \clearpage
    \thispagestyle{empty}
    \null
    \cleardoublepage
}

%create abstract environment that is not available in latex book style


% --- Basiskodierung und Fonts (optional, je nach Setup) ---
\usepackage[T1]{fontenc}
\usepackage[utf8]{inputenc}
\usepackage{lmodern}

% Palatino als Serif, Helvetica als Sans
\usepackage{palatino}
\usepackage{helvet}
\renewcommand{\sfdefault}{phv} % explizit Helvetica
\linespread{1.05}

% Mikrotypografie
\usepackage{microtype}

% Layout / Geometrie
\usepackage{geometry}
\geometry{a4paper, left=2.5cm, right=2.5cm, top=4cm,
          bottom=5cm, bindingoffset=1cm}

% Zeilenabstand
\usepackage{setspace}
\onehalfspacing

% Absätze
\usepackage{parskip}
\sloppy

% Überschriftenformatierung
\usepackage{titlesec}
\titleformat{\chapter}[display]
  {\LARGE\bfseries}{\chaptertitlename\ \thechapter}{10pt}{\huge\bfseries}

% Farben (für listings etc.)
\usepackage{xcolor}
\definecolor{darkblue}{rgb}{0.0,0.0,0.5}
\definecolor{grey}{rgb}{0.8,0.8,0.8}
\definecolor{lightgrey}{rgb}{0.95,0.95,0.95}

% Listings
\usepackage{listings}
\lstset{
  language=Python,
  basicstyle=\footnotesize\ttfamily,
  captionpos=b,
  frame=tb,
  commentstyle=\color{gray}\bfseries,
  stringstyle=\ttfamily,
  keywordstyle=\color{darkblue}\bfseries,
  breaklines=true,
  aboveskip=10mm,
  belowskip=10mm,
  showstringspaces=false, % ← hier war bei dir das Komma fehlend
  numbers=left,
  % stepnumber=5,
  numberstyle=\tiny,
  numbersep=5pt
}

% comment-Umgebung
\usepackage{comment}

% Bezeichnungen anpassen
\renewcommand{\bibname}{References}
\renewcommand{\lstlistlistingname}{List of Listings}

% Für \url und klickbare Links (falls du DOIs/URLs nutzt)
\usepackage{hyperref}



% --- Metadaten ---
\newcommand{\mytype}{Investigation and Prototypical Implementation of Automated Component Integrations in XNAT}
\newcommand{\mycourse}{Bachelor}
\newcommand{\myauthor}{Tanae Bousfiha}
\newcommand{\mydepartment}{HAWK Hochschule für angewandte wissenschaft und kunst 
\\Institut für Medizinische Informatik der Universitätsmedizin Göttingen}
\newcommand{\mysubmissiondate}{xx. August 2025}
\newcommand{\myfirstsupervisor}{Prof.Dr. Roman Grothausmann}

\begin{document}

  \pagenumbering{roman}
  \setcounter{page}{1}

  %--- cover page ---
  %
% title page
%

\begin{titlepage}
    %--- logo ---
    \normalsize
    \begin{tabularx}{\textwidth}{lXr}
        \multirow{2}{*}{\includegraphics[width=6.5cm]{images/goe-logo.jpg}} 
        & & \\
    \end{tabularx}

    %default settings for the rest
    \large
    \centering

    \vspace{3cm}

    \textbf{\LARGE \mytype}\\

    submitted in partial fulfillment of the\\
    requirements for the course ``\mycourse''

    \vspace{2cm}

    \textbf{\LARGE \ The Automation of implementation of Docker Container in XNAT}

    \vspace{2cm}

    \myauthor{ }

    \vspace{2cm}

    \mydepartment

    \vspace{2cm}

    Bachelor's and Master's Theses\\
    of the Center for Computational Sciences\\
    at the Georg-August-Universität Göttingen

    \vspace{0.2cm}

    \mysubmissiondate


    %--- new page ---
    \myemptypage
    \clearpage
    \thispagestyle{empty}
    \null
    \flushleft
    \onehalfspacing
    \normalsize

    \vspace{12cm}

    Georg-August-Universität Göttingen\\
    Institute of Computer Science\\[3ex]
    Goldschmidtstraße 7\\
    37077 Göttingen\\
    Germany\\[3ex]

    \begin{tabular}{@{}ll}
        \Telefon & +49 (551) 39-172000\\
        \fax & +49 (551) 39-14403\\
        \Letter & \href{mailto:office@informatik.uni-goettingen.de}{office@informatik.uni-goettingen.de}\\
        \Mundus & \url{www.informatik.uni-goettingen.de}\\
    \end{tabular}

    \vspace{1.0cm}

    \begin{tabular}{@{}ll}
        First Supervisor: & \myfirstsupervisor\\
        Second Supervisor:& \mysecondsupervisor\\
    \end{tabular}

    \clearpage
\end{titlepage}



  \section{Acknowledgments}
THANK YOUU

  %--- statement page ---
  
\includepdf[pages=-]{slbki.pdf}

  %--- abstract ---
  \clearpage\phantomsection\pdfbookmark{\abstractname}{abstract}
  \thispagestyle{empty}
  %
\begin{abstract}
This Work investigates the exiting methods of implementation in \ac{XNAT}(eXtensible Neuroimaging Archive Toolkit), with the aim of identifying limitations and best practices. Based on this analysis, an automated implementation workflow was developed. It was hypothesized that a Container could be implemented in a Project in XNAT and to rework all the files out and to upload the result files back to their places. The approach is partly based on using a Python Script that worked with \ac{REST API} and create a Dockerfile and asks the user for an external/ Container Script and extract all the files from the Open source XNAT. Upon the conclusion of the experience, the Docker Container in XNAT should be in ‘Complete’ turned and the result files should be uploaded. Experimental application of the methods demonstrated that the container could not receive the extracted files via REST API, which led to a not feasible neither processing nor uploading files. Results revealed a significant correlation between the workflow data of XNAT and the Mounting of data in the Docker container. These findings demonstrate that the extraction of all the files from all the levels of XNAT leads to a problem Mounting in the XNAT host before arriving to the container. And the fact that the Container has no access to data Bank of XNAT conducts that the effectiveness of the method proved to be context dependent and thus limited.
\end{abstract}



  % reset acronyms after abstract
  \acresetall

  %--- table of contents ---
  \clearpage\phantomsection\pdfbookmark{\contentsname}{toc}
  \tableofcontents

  %--- lists ---
  \listoffigures
  \listoftables
  \lstlistoflistings

  %--- list of acronyms ---
  




\chapter*{List of abbreviations}
\begin{acronym}
  \acro{XNAT}{Extensible Neuroimaging Archive Toolkit}
  \acro{OSA}{Obstructive Sleep Apnea}
  \acro{MRI}{Magnetic Resonance Imaging}
  \acro{AI}{Artificial Intelligence}
  \acro{ID}{Identifier}
  \acro{DICOM}{Digital Imaging and Communications in Medicine}
  \acro{ML}{Machine Learning}
  \acro{JSON}{JavaScript Object Notation}
  \acro{FAIR}{Findable, Accessible, Interoperable, Reusable}
  \acro{Stdout}{Standard Output}
  \acro{UI}{User Interface}
  \acro{HTTP}{Hypertext Transfer Protocol}
  \acro{XML}{Extensible Markup Language}
  \acro{REST}{Representational State Transfer}
  \acro{API}{Application Programming Interface}
\end{acronym}






  % arabic page numbers
  \pagenumbering{arabic}
  \setcounter{page}{1}

  %--- content ---
  

\chapter{Introduction}

Artificial intelligence holds significant, largely untapped potential within healthcare. Successful implementation of \ac{AI} could alleviate workload pressures on healthcare professionals while simultaneously enhancing the quality of their work through error reduction and increased precision. Furthermore, AI could empower patients to take a more active role in managing their own health and contribute to a reduction in avoidable hospital admissions~\cite{aung_promise_2021}.


XNAT, developed at Washington University’s Neuroinformatics Research Group, is an open-source platform for storing, managing, and analyzing imaging data. Its flexible design makes it suitable for many different imaging projects\footnote{\url{https://xnat.org/about/} (last accessed: 24.09.2025)}.


The open-source XNAT archiving system has expanded to include biosignal data and integrated machine learning, allowing AI models to transform it into an active analysis hub through automated processing~\cite{marcus_extensible_2007}.

In the Somnolink project, XNAT facilitates interoperable sleep data exchange to improve \ac{OSA} from diagnosis to therapy. AI systems integrated within XNAT support early patient identification, treatment planning, and therapy  adherence~\cite{internetredaktion_somnolink_nodate}.

Studies on AI tools in clinical research often emphasize particular technical infrastructures, neglecting generalizable implementation strategies within data platforms. For instance, by combining NVIDIA’s Clara Deploy with a customized XNAT, Hawkins et al. (2023)~\cite{hawkins_implementation_2023} showed how machine-learning pipelines can be integrated into clinical workflows for near real-time processing. This approach transformed XNAT into an active hub for pipeline-driven analysis, with results stored for clinician review. Building on this, Satrajit Chakrabarty et al. ~\cite{chakrabarty_deep_2023} presented an end-to-end framework capable of automating \ac{MRI} scan classification and tumor segmentation using AI-based methods in neuro-oncology studies.

Although these examples demonstrate technical feasibility, existing studies often emphasize specific infrastructures or disease context rather than examining broader, generalizable implementation strategies within data platforms. This thesis addresses that gap by investigating the methods for integrating \ac{ML} within XNAT, with the goal of identifying best practices, recognizing limitations, and proposing future measures to support scalable and generalizable AI implementations in healthcare data platforms.

This bachelor thesis will examine the Docker Container approach within XNAT, considering its benefits, disadvantages, and future development. It will assess its relevance to the Somnolink project, particularly regarding sustainability and FAIR principles. The findings will be synthesized into a clear User Guide for the practical application of the Docker Container approach.
The study will also analyze the XNAT Pipeline approach and compare it to the Docker Container method, exploring future measurement possibilities for both.

Chapter 2 lays the fundamentals by explaining the primary concepts necessary for understanding the thesis. Chapter 3 then details the Docker Container method, explaining its deployment within XNAT and introducing the automation script. A user guide is provided for users. Following this, the thesis addresses the Pipeline method and its deployment within XNAT.

In Chapter 4 provides a comprehensive analysis of both the Docker Container and Pipeline methods within XNAT. It examines their  generalizability, relevance to the Somnolink project including sustainability and FAIR criteria, and their advantages and disadvantages. Furthermore the chapter explores potential future measurements to further assess both methods.



 


  
\chapter{Fundamentals}

An overview of the system and methods used is provided in this chapter, serving as a foundation for the more detailed examination of the approaches that follows. 

\section{The Extensible Neuroimaging Archive Toolkit}
The Extensible Neuroimaging Archive Toolkit (XNAT)~\cite{marcus_extensible_2007} is an open-source imaging informatics platform developed by the Neuroinformatics Research Group at Washington University. It facilitates common management, productivity, and quality assurance tasks for imaging and associated data. Thanks to its extensibility, XNAT can be used to support a wide range of imaging-based projects.\footnote{\url{https://www.xnat.org/about/} (accessed on 10 August 2025)} Due to the fact that XNAT offers the opportunity to use the power of high-performance computing on your data~\cite{zaschke_extending_2024}, 
pipeline processing is one of the options that can be done on XNAT, as well as the deployment of Docker containers.
\\
Additionally, XNAT has a structured website that is composed of a Dashboard that gives an overview of the site, a Project list that browses or searches available projects, and a Project page where all subjects and their data in this project are located. Along with Subject/Session pages where all details are given, including imaging sessions and scan-level details, as well as the section of scans where all the patient’s scans are found. 
This structure of XNAT provides the information on where the patient files are stored, including the composed URL for REST API use. A REST API is a type of application programming interface that adheres to the principles of the REST architecture. REST, which stands for Representational State Transfer, defines a set of conventions for designing web APIs in a consistent and scalable way.\footnote{\url{https://www.redhat.com/en/topics/api/what-is-a-rest-api} (accessed on 10 August 2025)}
\\

\section{The REST API}
REST is defined by its architectural contrainst, not by specific protocol or standars, such as API, allowing for various implementation approaches by API developers.
An API is a set of terminology and standards for building and integrating application software. Those two are the main ingredients to achieve a automation of the analytical methods. \footnote{https://www.redhat.com/en/topics/api/what-is-a-rest-api{last acceed 07.09.2025}}

\section{The Container}
A container is a standard unit of software that packages up code and all its dependencies so the application runs quickly and reliably from one computing environment to another.\footnote{\url{https://www.docker.com/resources/what-container/} (accessed on 10 August 2025)}
Docker is a platform that allows certain applications to run in software containers. It was launched in 2013.\footnote{\url{https://fr.wikipedia.org/wiki/Docker\_(logiciel)} (accessed on 10 August 2025)} A Docker container allows the user to deploy any application in any environment without worrying about incompatibility issues, regardless of the machine’s configuration settings.\footnote{\url{https://sematext.com/glossary/docker/} (accessed on 10 August 2025)} Although there are some associated terms that will always be needed when it comes to using or learning about Docker containers.
One of the most important of them is the Dockerfile. The Dockerfile is a file that contains the build instructions for the image that the Docker Engine (an open-source containerization technology for building and containerizing your applications\footnote{\url{https://docs.docker.com/engine/} (accessed on 10 August 2025)}) will run. Conversely, there are numerous commands that must be written in the Dockerfile in order to create the image. After building the image, we can tag and push the image to Docker Hub. For this purpose, we just have to run some basic simple commands on a terminal.
\\
Another element to consider is the REST APIs. A REST API is an application programming interface (API) that follows the design principles of the REST architectural style. REST is short for Representational State Transfer and is a set of rules and guidelines about how you should build a web API.\footnote{\url{https://www.redhat.com/en/topics/api/what-is-a-rest-api} (accessed on 10 August 2025)} Luckily, XNAT provides its users with a list of REST APIs that make the exchange between the user and the website possible. 

\section{Command Definition in JSON format}
At the same time, if we want XNAT to execute a created Docker image and run a command, a \ac{JSON} command definition in JSON format is needed. 

The command definition in JSON format is a structured set of properties that defines the Docker image, which enables XNAT to execute the command correctly.\footnote{\url{https://wiki.xnat.org/container-service/json-command-definition}, last accessed 23.08.2025}

In other words, this command definition in JSON format (written in JSON) can be seen as a configuration for XNAT. When defining a command for XNAT in JSON format, the Docker image and a clear, human-readable name are specified, along with the exact command-line invocation, what input files are required and where they are mounted, how those inputs are selected from XNAT, and what outputs are produced and where they should be uploaded back into XNAT.\footnote{\url{https://wiki.xnat.org/container-service/json-command-definition}, accessed 23 August 2025}

\section{The Pipeline}
Pipelines in XNAT are efficient lightweight applications that can be run on the project data to facilitate sophisticated processing tasks or leverage the computational power of extensive computing clusters. Some pipeline-enabled workflows run automatically, such as auto-QC of images, and some of them require manual steps, like defining a region of interest. After integrating the pipeline in project data, we can set up project-based workflows with project-specific and experiment-specific parameters, track a pipeline and send email notifications, and capture provenance information as the pipeline executes.\footnote{https://wiki.xnat.org/documentation/running-pipelines-in-xnat {last accessed 07.09.2025}} The pipelines are used for tasks like converting DICOM to NIFTI, performing image quality assessment (QA), correcting MRI bias fields, and running advanced analyses such as the FreeSurfer pipeline.\footnote{https://github.com/jhuguetn/xnat-pipelines {last accessed 07.09.2025}}

\section{The Pipeline Descriptor}

Pipelines are identified by using a pipeline descriptor and resource descriptors. A resource descriptor describes an executable. An executable is identified by its name, its location, and the arguments that it takes.\footnote{https://wiki.xnat.org/documentation/xnat-pipeline-development-schema {last accessed 07.09.2025}} In other words, the Pipeline Descriptor is a recipe of the pipeline that tells XNAT exactly what a pipeline is, what input is needed, and where the output should be located, ensuring a seamless integration in XNAT.









  \chapter{The Analysis of the First Method: Docker Container}

\section{Introduction to the Automation Script for the Docker Container Method}
The first method analyzed was the Docker Container, and after reviewed the last points supported by XNAT Community, a review of the latest updates from the XNAT Community indicated that automation of Docker container implementation is feasible.

\section{Table of the REST API in the Script}

In order to summarize all the REST API needed in the script, a summarizing table was created.

\begin{longtable}{|l|p{0.4\textwidth}|p{0.4\textwidth}|}
    \hline
    \textbf{Method} & \textbf{Endpoint} & \textbf{Purpose} \\
    \hline
    \endfirsthead

    \hline
    \textbf{Method} & \textbf{Endpoint} & \textbf{Purpose} \\
    \hline
    \endhead

    POST & \path{/xapi/commands} & Upload a Command (JSON) \\
    \hline
    GET  & \path{/xapi/commands} & List all the commands to get the wrapper ID and the command ID \\
    \hline
    PUT  & \path{/xapi/commands/{command_id}/wrappers/{wrapper}/enabled} & Enable wrapper site-wide \\
    \hline
    POST & \path{/xapi/projects/{project_id}/commands/{command_id}/wrappers/{wrapper}/root/project/launch} & Launch container wrapper in project context \\
    \hline
    GET  & \path{/data/projects/{project_id}/resources?format=json} & List all project resources \\
    \hline
    GET  & \path{/data/projects/{project_id}/resources/{resource_label}/files?format=json} & List files inside a project resource \\
    \hline
    GET  & \path{/data/projects/{project_id}/subjects?format=json} & List all subjects in a project \\
    \hline
    GET  & \path{/data/subjects/{subject_id}/resources?format=json} & List subject-level resources \\
    \hline
    GET  & \path{/data/subjects/{subject_id}/resources/{resource_label}/files?format=json} & List files inside a subject resource \\
    \hline
         & \path{/data/subjects/{subject_id}/resources/{resource_label}/files/{file_name}} & Direct file path (download or reference) \\
    \hline
    GET  & \path{/data/projects/{project_id}/experiments?format=json} & List sessions (experiments) in a project \\
    \hline
    GET  & \path{/data/experiments/{experiment_id}/resources?format=json} & List session-level resources \\
    \hline
    GET  & \path{/data/experiments/{experiment_id}/resources/{resource_label}/files?format=json} & List files in a session resource \\
    \hline
         & \path{/data/experiments/{experiment_id}/resources/{resource_label}/files/{file_name}} & Direct file path (download or reference) \\
    \hline
    GET  & \path{/data/experiments/{experiment_id}/scans?format=json} & List scans within a session \\
    \hline
    GET  & \path{/data/experiments/{experiment_id}/scans/{scan_id}/resources?format=json} & List scan-level resources \\
    \hline
    GET  & \path{/data/experiments/{experiment_id}/scans/{scan_id}/resources/{resource_label}/files?format=json} & List files inside a scan resource \\
    \hline
      
\end{longtable}


The REST API is a sort of backbone of the automatization in XNAT workflow, it assures a orchestration of the container execution, and at the same time assuring  traceability and reproducibility. The REST API played a crucial role to develop a automatic workflow in the XNAT platform. The figure 3.1 illustrates the role of the REST API in the automatization


\begin{figure}[ht]
    \centering
    \def\svgwidth{0.9\linewidth}
    \input{RESTAPI.pdf_tex}
    \caption{Diagram: illustrates the role of the REST API in the automatization}
    \label{fig:workflow-steps}
\end{figure}














\section{The Automation Script}
The purpose of this chapter is to introduce the system in more detail. A Python script is used to combine all the steps into an automated workflow. The libraries used in this work include: \texttt{requests}\footnote{\url{https://www.w3schools.com/python/module_requests.asp} accessed on 10 August 2025}, \texttt{json}\footnote{\url{https://www.w3schools.com/python/python_json.asp} accessed on 10 August 2025}, \texttt{os}\footnote{\url{https://docs.python.org/3/library/os.html} accessed on 10 August 2025}, \texttt{subprocess}\footnote{\url{https://docs.python.org/3/library/subprocess.html} accessed on 10 August 2025}, \texttt{getpass}\footnote{\url{https://docs.python.org/3/library/getpass.html} accessed on 10 August 2025}, \texttt{sys}\footnote{\url{https://docs.python.org/3/library/sys.html} accessed on 10 August 2025}, and \texttt{urllib3}\footnote{\url{https://urllib3.readthedocs.io/en/stable/} accessed on 10 August 2025}.

The use of the Requests library was for creating the interaction with the web services of XNAT. The Python Requests library is the go-to solution for making HTTP requests in Python, thanks to its elegant and intuitive API that simplifies the process of interacting with web services and consuming data in the application.\footnote{\url{https://www.datacamp.com/tutorial/python-subprocess} accessed on 10 August 2025} The Sys module in Python provides access to variables and functions that interact closely with the Python interpreter and runtime environment.

Since a JSON command is written, the use of the JSON library was necessary. JSON is a syntax for storing and exchanging data. It is text, written with JavaScript Object Notation.\footnote{W3Schools, Python JSON}
The Python Subprocess module is a tool that allows running other programs or commands from your Python code. Using the Python Subprocess module resembles executing commands directly on the system computer using Python instead of typing them directly into the command prompt. This module makes it easy to automate Python code.\footnote{\url{https://www.datacamp.com/tutorial/python-subprocess} accessed on 10 August 2025}
Getpass() prompts the user for a password without echoing. The getpass module provides a secure way to handle password prompts where programs interact with the users via the terminal.\footnote{\url{https://www.geeksforgeeks.org/python/getpass-and-getuser-in-python-password-without-echo/} accessed on 10 August 2025}

The urllib3 library was used to disable the InsecureRequestWarning, which is normally shown when Requests connects to a site with an unverified SSL certificate. This approach helps keep the output clean and free of unnecessary warnings during the run of the script.

Currently, the system is running Docker Engine 28.3.0. The Docker Client version in use is 20.10.5+dfsg1, which communicates with the Docker Engine to issue commands and manage images and containers. The XNAT web application is running on version 1.9.1.1, as verified through the web interface.


\begin{figure}[ht]
  \centering
  \def\svgwidth{0.8\linewidth}
  \input{lib.pdf_tex}
  \caption{Schema: Core Libraries for Python Scripting in System Integration}
  \label{fig:diagram-core-libraries}
\end{figure}





\section{Architecture and organization of the Code}
The automation script is structured into 17 distinct parts, each one responsible for a specific step in the overall workflow.
During the development of the automation script, I closely followed the steps of the manual process, since I had already tested and completed it manually. This included writing a container script, writing a Dockerfile, building the image, writing a JSON command, and enabling the command first on the website and then on the project, with the last step being the launch of the container. I carefully repeated the same steps during the implementation. In other words, every structured part is composed of a function that is responsible for carrying out a step of the manual process.
 
 \section{The Dockerfile}
 
The first function responsible for writing the Dockerfile is called \texttt{def write\_dockerfile}, and it takes three arguments: the target directory for the Dockerfile, the name of the Python script to include, and an optional base Docker image (defaulting to \texttt{python:3.10-slim}, currently the fastest and slimmest image of Docker). The function creates the necessary folder (if it does not already exist), constructs the Dockerfile content, and writes it. The content of the Dockerfile is introduced with string interpolation and contains the base image and several commands such as \texttt{WORKDIR /app} (a working directory) and \texttt{COPY \{script\_filename\} /app/\{script\_filename\}} to copy the user script inside the image. The line \texttt{RUN pip install --no-cache-dir pandas} tells Docker to install the pandas Python library using pip during the build process of the Docker image, which helps reduce the overall image size and avoids unnecessary layers in the build process. The last component of the Dockerfile is \texttt{COPY requirements.txt /app/requirements.txt}, which provides the opportunity for the user to write in a separate text file all the libraries used in the container script, to be installed later during the image build process.


The \texttt{CMD} instruction was removed from the Dockerfile because it caused errors when running the container in XNAT. The issue was due to a conflict between the \texttt{CMD} instruction in the Dockerfile and the JSON command configuration, as both included the \texttt{python3} prefix. This led to a duplicate command execution, which caused the container to fail. By removing \texttt{CMD} from the Dockerfile, the execution logic is now fully controlled by XNAT through the JSON command, ensuring a successful launch of the container.

\begin{lstlisting}[language=bash,caption={Dockerfile}]
FROM python:3.10-slim
WORKDIR /app
COPY script.py /app/script.py
RUN pip install --no-cache-dir pandas
COPY requirements.txt /app/requirements.txt
\end{lstlisting}

The rest of the \texttt{def write\_dockerfile} ensures that the specified directory exists, then creates and writes a Dockerfile to that location. It constructs the full path to the file, writes the generated content into it, and prints a confirmation message.


\lstset{
  language=Python,
  basicstyle=\ttfamily\small\color{black},
  keywordstyle=\color{black},
  identifierstyle=\color{black},
  stringstyle=\color{black},
  commentstyle=\color{black},
  numberstyle=\color{black},
  showstringspaces=false
}

\begin{lstlisting}
os.makedirs(docker_dir, exist_ok=True)
dockerfile_path = os.path.join(docker_dir, "Dockerfile")
with open(dockerfile_path, "w") as f:
    f.write(dockerfile_content)
print(f"Dockerfile written to {dockerfile_path}")
return dockerfile_path
\end{lstlisting}

This function ensures that the Dockerfile is written in an appropriate way and secures that all the dependencies are stored in the external requirements text file. It handles all possible cases and avoids the issues of errors and standard input errors.

\section{Building, Pushing and Tagging the Image}

When a command is deployed in XNAT and a container is launched, the first thing that XNAT checks is whether the Docker image is available in DockerHub. This means the user must have an account in DockerHub and preferably log in before using the following script. Consequently, it is essential to push and tag the created image.

The function responsible for the build is \texttt{def build\_docker\_image}, and it expects two parameters: \texttt{dockerfile\_path} for the Dockerfile path and the name/tag for the Docker image to be built \texttt{docker\_image\_name}. For the build, we used \texttt{subprocess.run} to call the external \texttt{docker build} command.


Usually, when we want to build an image in Docker, we use the basic command \texttt{docker build .}

\begin{lstlisting}
full_tag = f"{dockerhub_username}/{docker_image_name}"
    print(f"Building Docker image '{full_tag}'...")
    build_result = subprocess.run(
        ["docker", "build", "-f", dockerfile_path, "-t", full_tag, "."],
        capture_output=True, text=True)
    if build_result.returncode != 0:
        print(f"Build failed:\n{build_result.stderr}")
        sys.exit(1)
    print(f"Image '{full_tag}' built successfully.")
\end{lstlisting}
As we can see in this block, \texttt{subprocess.run} was used to call the external \texttt{docker build} command, and the \texttt{-f dockerfile\_path} option tells Docker which Dockerfile to use. The \texttt{-t docker\_image\_name} option sets the name and tag for the image. The point (.) at the end of the command means the current directory will be used as the build context. In addition, \texttt{capture\_output=True} collects both stdout and stderr so the script can handle their output or errors, and \texttt{text=True} ensures that outputs are returned as strings.

Now that the image has been successfully built and tagged, we can proceed to the next essential step: pushing the image to Docker Hub. To do this, the same method of \texttt{subprocess.run} is used.


\begin{lstlisting}
    push_result = subprocess.run(["docker", "push", full_tag], capture_output=True, text=True)
    if push_result.returncode != 0:
        print(f"Push failed:\n{push_result.stderr}")
        sys.exit(1)

    print(f"Image successfully pushed: {full_tag}")
    return full_tag
\end{lstlisting}

In normal cases, we use the command \texttt{docker push image\_name}, and this assumes that the user is already logged in to their Docker account. However, in the automation script, we make use of \texttt{subprocess.run}.

To explain this block, we begin by discussing the reason for using the full image tag. It constructs a complete image tag in the Docker Hub format. Since Docker Hub requires image tags to follow a specific structure, the full image tag uses the format \texttt{dockerhub\_username/image\_name[:tag]}.

This line creates a correctly formatted name for the Docker image, linking the Docker Hub account \texttt{(dockerhub\_username)} with the chosen image name and optional version tag \texttt{(docker\_image\_name)}.

The line of the code: 
\begin{lstlisting}
push_result = subprocess.run(["docker", "push", full_tag], capture_output=True, text=True)
\end{lstlisting}
refers to the push process with the \texttt{subprocess.run} method. And it follows the same concept as building the Docker image.

\section{The Prompt Function for the Required Input}

This function performs the task of capturing input from the user.
We have used this function in the script in multiple scenarios. The first was to take the Docker Hub username from the user. To customize the JSON command, we asked the user for the name of the command and its description. Overall, it is also used to take the username and password from the user’s login credentials. The function proceeds through all of these steps while running the script.

The function looks like this:
 
 \begin{lstlisting}
 def get_input(prompt):
    while True:
        value = input(prompt)
        if value.strip():
            return value
        else:
            print("Cannot be empty.")

\end{lstlisting}
The \texttt{def get\_input} function uses a \texttt{while True} loop to keep asking until the user enters input. The \texttt{value.strip()} removes any whitespace from the user input. If the input is valid, the value is returned; if not, the message \texttt{Cannot be empty} will be printed.

\section{The JSON Command}

To communicate the Docker image to XNAT, and consequently to run the container in the website, it is necessary to write a JSON command.\footnote{\url{https://wiki.xnat.org/container-service/json-command-definition}}
Usually, the JSON command is personalized and customized depending on the main idea of the image. But in the automation case, the JSON command must be generalized for all cases, which means no matter which file is selected, the command will still handle it.

To achieve that, the function used is called \texttt{create\_json\_file}, which builds the configuration dictionary for a Docker-based XNAT command (for the XNAT Container Service), writes it to a \texttt{command.json} file, and returns the filename.

The function expects three parameters:

\texttt{docker\_image}: The Docker image to use (string).

\texttt{script\_filename}: The name of the Python script inside the container (string).

\texttt{mod\_data}: A dictionary holding user-provided metadata (names, descriptions, etc.).

Let us start analyzing the first block of the JSON command:
\begin{lstlisting}

json_file = {
        "name": mod_data["command_name"],
        "description": mod_data["command_description"],
        "version": "1.5",
        "type": "docker",
        "image": docker_image,
        "command-line": f"python3 /app/{script_filename} /input/#input_file# /output",
        "mounts": [
            {"name": "input", "path": "/input", "writable": False},
            {"name": "output", "path": "/output", "writable": True}
        ],
\end{lstlisting}


In this block, we are defining the name of the command, adding a description, and specifying the version and the type of the image.
The command line declares the actual command to run inside the container when invoked by XNAT. The placeholder \# is a template substitution (not a regex expression) that tells XNAT: \texttt{"When you launch the Docker command, replace \#input\_file\# with the actual file name/path that the user selected as input."} It does not define what a valid file looks like; instead, it marks a spot in the command where XNAT should "fill in" an actual value.


The mount configures the directory mappings (as Docker volumes) between XNAT and the Docker container. Most of the time, two are used: the input and the output.
We specified the path \texttt{/input} in the container and indicated whether the container can only read the files (\texttt{Not writable}) or can also write them (\texttt{Writable}).

The second part of the JSON command:


\begin{lstlisting}

 "inputs": [
            {
                "name": "input_files",
                "description": "Input files",
                "type": "files",
                "element_type": "file",
                "required": True,
                "mount": "input",
                "multiple": True
            }
        ],
        "outputs": [
            {
                "name": "result_file",
                "type": "file",
                "description": "Result file output",
                "mount": "output",
                "path": "."
            }
        ],
\end{lstlisting}


This block defines the parameters that a user must provide as an input file when launching the container.
Let’s break down each field:

The \texttt{"name": "input\_files"} is used in other parts of the JSON as a reference---such as placeholders in the command line. In addition, we add an optional human-friendly description. The \texttt{"type"} indicates that the input accepts more than just text files, scans, or numbers; it specifies that each element in the input is specifically a file through the \texttt{"element"} field.

The \texttt{"required": true} flag means that if the input is not provided, the command cannot be run. The mount in this part links the input to a specific mount inside the container. All files provided by the user will appear in the \texttt{/input} directory inside the container. With \texttt{"multiple": true}, it is indicated that the user can upload or select more than one file for this input.

The same logic applies to the output part. The only additional point to note is the \texttt{"path": "."}, which means the output file will be placed at the top level of the \texttt{/output} directory. To allow more than one output file, the setting \texttt{"type": "files"} is used in the output block.

The final block in the JSON command:

\begin{lstlisting}

 "xnat": [
            {
                "name": wrapper_name,
                "label": mod_data["label_name"],
                "description": mod_data["label_description"],
                "contexts": ["xnat:projectData"],  
                "external-inputs": [
                    { 
                        "name": "project",
                        "type": "Project",
                        "required": True,
                        "load-children": True
                    }
                ],
                "output-handlers": [
                    {
                        "name": "output",
                        "accepts-command-output": "result_file",
                        "as-a-child-of": "project",
                        "type": "Resource",
                        "label": "Results",
                    }
\end{lstlisting}

This block defines what will be displayed on the XNAT interface and connects the Docker command with XNAT's web interface and permission system. It controls how the command appears as a tool in XNAT, how it is integrated with projects or other data, and how inputs are handled.

The wrapper name is the technical name of the command inside XNAT. The \texttt{"label": mod\_data["label\_name"]} is a human-readable name shown in the XNAT \ac{UI}. The \texttt{"description": mod\_data["label\_description"]} is a description for the tool/wrapper, shown when users browse tools in XNAT. The \texttt{"context"} specifies where we want the container to appear in the XNAT interface. Since the main idea was to have one container on top of the structure of XNAT, here \texttt{["xnat:projectData"]} means this command is available at the project level.


In the \texttt{"external-input"} block, external entities from XNAT are defined. Because we used the project as the context for the command, we require in the external input that a project be selected. In more detail, the name and type must be "Project." This input must be provided, which is why \texttt{"required": true} is specified. We tell XNAT with \texttt{"load-children": true} to load (in the UI) the child objects of the project when showing this input.

In the \texttt{"output-handlers"} block, we control how the outputs from the command are stored and shown in XNAT after job completion. This output handler tells XNAT to take the result file produced by the command, store it as a resource under the project, and call that section ``Results'' in the UI. The \texttt{"as-a-child-of": "project"} specifies that the results are uploaded back into the XNAT project as new resources.

Finally, the script writes the JSON command file, ensuring that the JSON command is saved in the correct format, ready to be uploaded to XNAT.

\section{Upload the Command in XNAT}

After writing the JSON command, the next step is to send it to XNAT. To achieve this, we have to use the appropriate REST API responsible for uploading the container command to XNAT.
We can find the list of all REST APIs under \texttt{Administer $\rightarrow$ Site Administration $\rightarrow$ Miscellaneous $\rightarrow$ Development Utilities $\rightarrow$ Swagger}.
To send the command to XNAT, we found that the responsible REST API is \texttt{POST}.\footnote{\url{https://xnat-dev.gwdg.de/xapi/swagger-ui.html\#/command-rest-api}}

The function responsible for this is \texttt{def send\_json\_to\_xnat}, and it expects four parameters: \texttt{json\_file\_path, xnat\_url, xnat\_user, xnat\_password}. When working with REST APIs, we first have to build the URL endpoint for XNAT’s command registration.
 
\begin{lstlisting}
def send_json_to_xnat(json_file_path, xnat_url, xnat_user, xnat_password): 

    url = f"{xnat_url}/xapi/commands"
    print(f"Uploading command to {url}")
    with open(json_file_path, "r") as f:
        response = requests.post(url, auth=(xnat_user, xnat_password), json=json.load(f))
    if response.status_code == 200:
        print("Command uploaded successfully.")
    elif response.status_code == 201:
        print("Command created successfully.")
    elif response.status_code == 409:
        print("Command already exists.")
    else:
        print(f"Failed to upload command: {response.status_code} - {response.text}")
\end{lstlisting}

This function takes the path that describes the JSON command, the base URL of the XNAT server, and the credentials for authentication. First, it constructs the correct REST API endpoint by appending \texttt{/xapi/commands} to the provided XNAT URL, ensuring that the command is sent to the appropriate API for command registration. It then opens the JSON file, loads its contents, and sends this data as a POST request to the API endpoint using \ac{HTTP} Basic Authentication with the supplied username and password. After making the request, the function examines the server’s response, and in each case provides feedback, helping the user quickly understand whether the command upload succeeded or if further action is needed to resolve an issue.

\section{Preparations to Launch the Container}

The goal of the previous steps is to launch the container in XNAT, but before that we need to perform some preparations. In line with the automation principle, we need to retrieve some information from the XNAT web service, and in this case we use the REST API technique as well. The first information that we need to extract after sending the JSON command to XNAT is the \texttt{Command ID} and the \texttt{Wrapper ID}.
In order to do that...
 
\begin{lstlisting}
def get_command_wrapper_id(xnat_host, xnat_user, xnat_password, command_name, wrapper_name=None):
 
    url = f"{xnat_host.rstrip('/')}/xapi/commands"
    try:
        resp = requests.get(url, auth=(xnat_user, xnat_password), verify=False)
    except Exception as e:
        print(f"Verbindungsfehler: {e}")
        sys.exit(1)
    if resp.status_code != 200:
        print(f"Fehler beim Abrufen der Commands: {resp.status_code}")
        sys.exit(1)
    data = resp.json()
    commands = data.get("commands", data) if isinstance(data, dict) else data

    for command in commands:
        if command.get("name") == command_name:
            if not wrapper_name:
                return command.get("id")
            
            for wrappers_field in ["xnat", "wrappers"]:
                for wrapper in command.get(wrappers_field, []):
                    if wrapper.get("name") == wrapper_name:
                        return wrapper.get("id") or wrapper_name
            print("no Wrapper found")
            sys.exit(1)
    print("Command nicht gefunden.")
    sys.exit(1)
\end{lstlisting}

The function takes as input the XNAT server’s hostname, user password, the name of the desired command, and the name of a specific wrapper. It first builds the correct XNAT REST API endpoint to list all the registered commands on XNAT.
Upon receiving a response, the function checks for a successful status (200 OK); otherwise, it reports the failure and exits. If only a command ID is needed (no wrapper specified), it returns the command’s ID directly. If a wrapper name is provided, the function searches for the wrapper within possible sections of the command definition under XNAT or the \texttt{"wrappers"} fields—returning the wrapper’s ID upon match. If the specified command or wrapper cannot be found, the function clearly reports this to the user.

After receiving the Command ID and the Wrapper ID, we proceed by enabling the command in the project and site-wide. For that, we use the two functions \texttt{def enable\_wrapper\_sitewide} and \texttt{def enable\_wrapper\_for\_project}, which both expect parameters such as the server’s hostname, user login and password, and the extracted Command ID and Wrapper ID. Each function builds the desired XNAT REST API endpoint (in this case, we used the \texttt{PUT} REST API), checks the status of the process, and provides feedback to the user to track the code’s progress.

\section{The Extraction of All Files from All Project Structures}

The function used in this part is \texttt{def get\_all\_files\_all\_levels}, and it collects all the files in the project data hierarchy (project, subjects, experiments/sessions, and scans) from an XNAT system using its REST API. We used the REST API endpoint \texttt{GET} for the extraction. The function proceeds by authenticating with the XNAT server using the provided user credentials and starts by gathering project-level files, querying all resource folders at the project root and listing every file found within them. At every step, the function structures each file into a dictionary containing hierarchical information (such as the level—project, subject, experiment, scan—and the relevant IDs), the file name, the resource folder it was found in, and an absolute URI for straightforward retrieval. By the end, it returns a comprehensive list of all file records across all levels in the project.
\begin{figure}[ht]
    \centering
    \def\svgwidth{0.9\linewidth}
    \input{extraction.pdf_tex}
    \caption{Diagram: The Extraction of files from all the project structure}
    \label{fig:enter-label}
\end{figure}

\section{Launch the Container}

Arriving at the well-known part of launching the container, we use the function \texttt{def launch\_container\_with\_all\_files}. This function is designed to automatically launch a container with all the files previously extracted as input. It takes the XNAT server connection details, project and command identifiers, user authentication, the wrapper name, and a list of files to be processed.

The function first checks if any files are provided. A payload dictionary is then prepared, mapping the project ID and the input files string to the fields. Next, the function constructs the appropriate URL for launching the command. Before sending the launch request, the function prints diagnostic information, including the chosen files and payload. It then submits a \texttt{POST} request with the payload as JSON and user credentials for authentication. After sending the request, the function reports back to the user about the status and the response.

\begin{lstlisting}
def launch_container_with_all_files(xnat_host, project_id, command_id, wrapper_name, xnat_user, xnat_password, files):
   
    if not files:
        print("no files found.")
        return 

    payload = {
        "project": project_id,
        "input_files": input_files_str
    }

    url = f"{xnat_host}/xapi/projects/{project_id}/commands/{command_id}/wrappers/{wrapper_name}/root/project/launch"

    headers = {"Content-Type": "application/json"}
    response = requests.post(
        url, auth=(xnat_user, xnat_password),
        headers=headers, json=payload, verify=False
    )

    print("Status:", response.status_code)
    print("Antwort:", response.text)
    if response.status_code in [200, 201, 202]:
        print("Container launched with all files!")
    else:
        print("failure :", response.status_code, response.text)
\end{lstlisting}

\begin{figure}[ht]
    \centering
    \def\svgwidth{0.9\linewidth}
    \input{steps.pdf_tex}
    \caption{Diagram: Workflow Steps for the Implementation}
    \label{fig:workflow-steps}
\end{figure}

The figure illustrates how the automation script facilitates the process of deploying a container in XNAT. It reduces all the steps.

\begin{figure}[ht]
    \centering
    \def\svgwidth{0.9\linewidth}
    \input{compare.pdf_tex}
    \caption{Diagram: Diagram illustrating the difference between the manual and the automatic process}
    \label{fig:enter-label}
\end{figure}

\begin{comment}
    https://wiki.xnat.org/documentation/strategies-for-xnat-image-data-storage
\end{comment}


\section{Testing the Code}
After testing the code, we can see in the beginning \texttt{see the Appendix at the end}, the process actively requests input from the user in order to proceed with the implementation. It started by asking the user for login credentials, then the Project ID that the user wants to process, and lastly, before proceeding with the creation and command sending to XNAT, it asked the user about the command name and the description.
After receiving all the information, the script commenced building the image, tagging, pushing, writing the JSON command, enabling the command, and lastly launching the container with all files.

\section{Results on XNAT}
The result that we came to after testing the code is that all the steps are working correctly, except that the container could not receive any files. With the \ac{STDout} view we found out that the container received zero files. 
 
\begin{lstlisting}[numbers=none]
View stdout (from file)
Inhalt von /input (rekursiv):
/input: []
Keine CSV-Datei in /input gefunden.
\end{lstlisting}


But interestingly, after searching in the container information, we noticed that the container in fact received the files and that the files are listed in the input part of the container input. 
Details of the test procedure are described in Appendix A \cite{bousfiha2025appendix}.

The results demonstrate two things. First, the automation script with the REST API worked correctly. Second, there is an issue with the container accessing the input files, despite them being listed as received.

\section{Discussion of the results}
This chapter presents an overview of the implementation solution, the technical challenges encountered during integration with XNAT, and the methods employed to overcome the issues. Several problems emerged throughout development, including data transmission failures, REST API inconsistencies, and JSON configuration mismatches.

The initial problem was that the container did not appear to receive any input files. To confirm this, a test was performed by adding a script within the container. The script was designed to write a file named \texttt{no\_file.txt} if no input file was detected at runtime. The successful upload of this file confirmed that while the container was running, it was not receiving actual input data.

Following the initial diagnosis, another test was conducted. Efforts were made to use the integration of the REST API within the container to retrieve files, process them, and upload the results. Despite the API being correctly called, no file results appeared in XNAT. 

Attention then shifted to the JSON command that defines the container command structure in XNAT. The JSON command was enhanced to include login credentials and more detailed output specifications. Multiple versions were tested, each with different contexts. Despite these efforts, no successful uploads occurred, and in some cases, the command itself was rejected by XNAT. 

It was hypothesized that the container execution timing could play a role in the failure. To test this, the automation script was modified to introduce a 10-second delay between each container execution. The assumption was that the host system might be under resource pressure when running multiple containers in rapid succession. However, introducing delays did not resolve the core issue: the container still did not receive input files.
%-------------------------------------------------------------------------------------------------------------------------



\section{Analysis of the Docker Approach: Current Limitations}

After analyzing the Docker Container method in both manual and automated scenarios, we identified several limitations, that requires further attention.
An advantage of using the Docker Container is the ability to automate the implementation workflow and launch either multiple containers or single workflows via scripting.

The Container approach was designed to be applicable on multiple structure of the XNAT. In other word depending on the intended case, users muss configure the container to process files within a specific structure of XNAT. However if the objective as to proceed all the files from all the existed structures in XNAT within a project, additional steps are needed, specifically extracting all the data from all the structures as well as extra work to relocate the processed files to their original directories after processing. 

\section{The Workflow Data in XNAT}

In most cases, when we press the button \textit{Run Container} in XNAT manually, the data is automatically provided, and the container handles the data and reloads the results back into the system.  
In detail, when the container is launched, the platform orchestrates a series of data management and processing steps to facilitate the workflow. At first, the system generates a JSON specification. This specification includes information about the script and command line to be executed, the output file location, and the container input, ensuring that the containerized process has access to all relevant instructions.  
Furthermore, XNAT identifies the relevant data and prepares it for transfer to the designated computational environment. The data is then mounted or copied to a temporary directory within the host system, ensuring accessibility for the containerized analysis. Once the container is activated, it processes the input data according to the command-line script, generates the output data, and finally writes the output back to the designated result directory on the host.  

Assuming that the container has the ability to process a large number of files, there is also the possibility to increase the number of files that can be handled using the command \texttt{ulimit}.  
The fact that XNAT mounts or copies the data to a temporary directory within the host system reveals important insights into the underlying problem. Copying the data to the host system and mounting it into a volume means that the container has no direct access to the XNAT database, which can potentially lead to issues related to memory or disk space on the host. This approach may result in resource constraints.  
This issue can become particularly problematic when large datasets are involved. Although my specific test case involved a relatively small number of files, the container still failed to process the data successfully. This suggests that the problem may not be solely related to dataset size or resource limits.

The diagram below illustrates the workflow data for a clearer understanding.  

\begin{figure}[ht]
    \centering
    \def\svgwidth{\linewidth} 
    \input{xnat.pdf_tex}
    \caption{Diagram: XNAT Workflow for Container Integration}
    \label{fig:workflowxnat}
\end{figure}


To recapitulate, the Docker container method proved effective for launching containers within multiple project structures. However, it did not perform well when attempting to process all files across a single project structure. This limitation highlights a clear restriction and reduces the overall usefulness of the approach.

\section{Generalizability of the Docker Automation Method}

Currently the Docker Container is a sort of powerful method to implement AI in the XNAT. it perfoms well on single models and also bound to a specific data type such as CSV,DICOM. However, adapting the workflow to accommodate other data types still requires manual intervention and adjustment. This limitation affects the method’s reusability in varying contexts or projects.
Generalizability is particularly important in a way that the method should be accessible to users without programming expertise and adaptable for use in different institutions, datasets, and environments.
One way to enhance the generalizability is through the implementation of a graphical user Interface (GUI). The GUI method can make the Docker Container method more accessible and understandable to all users even without programming experiences. Consequently facilitating
the adoption of this method in different environments, institutions and datasets.  
\section{GUI Mockup for Docker Container Automation in XNAT}

For the purpose of creating more accessibility and to generalizability for non technical or experienced users. The Docker automation script developed in this project was visualized as a graphical user interface (GUI) mockup.
The following mockup does not represent a functional program , but it illustrates how the command line workflow can be illustrated. In other word it a representation of what the result on the Terminal looks like. Consequently it follows the same logical steps as the one on the terminal: (1) configuration of the container, (2)project selection, (3)file discovery,  (4)container launch, and (4) result visualization.

\begin{figure}[ht]
  \centering
  \def\svgwidth{\linewidth}
  \import{./}{mockup1.pdf_tex}
  \caption{Diagram: The Mockup of a proposed interface for the Docker container automation in XNAT}
  \label{fig:workflowxnat}
\end{figure}

The figure  shows the proposed interface. 
  \chapter{The Analysis of the Second Method: The Pipeline}
\section{Pipeline Installation and Configuration in XNAT}

Pipeline is a method of implementation for automated data processing in XNAT, using \ac{XML} descriptors and server-side scripts/tools. 
To implement a pipeline in XNAT, an XML descriptor is required, this descriptor serves as a connector between XNAT and the user script including definitions for metadata, execution steps, conditions, iterations, and outputs—and resource descriptors that specify how scripts or executable are launched with necessary parameters.
During the Pipeline analysis we initially focused on analyzing the implementation process of a OSA prediction based script with the Pipeline method. The Pipeline system requires two components: Pipeline Engine and the XNAT Pipeline Engine Plugin. By default, pipelines are located at /data/home/xnat/pipeline\footnote{\url{https://wiki.xnat.org/documentation/xnat-pipeline-development-schema} last Accessed: 10 September 2025).}   

\begin{lstlisting}[language=XML, caption={Example of a pipeline resource descriptor in XNAT}, label={lst:pipeline-xml}]
<?xml version="1.0" encoding="UTF-8"?>
<pip:Resource xmlns:pip="http://nrg.wustl.edu/pipeline">
    <pip:name>run_test_pipeline.py</pip:name>
    <pip:location>/usr/local/python_scripts</pip:location>
    <pip:commandPrefix>python</pip:commandPrefix>
    <pip:type>Executable</pip:type>
    <pip:description>Generates SOMETHING HERE</pip:description>
    <pip:input>
        <pip:argument id="configfilepath">
            <pip:description>Path to param file</pip:description>
        </pip:argument>
        <pip:argument id="host">
            <pip:description>XNAT Host</pip:description>
        </pip:argument>
        <pip:argument id="user">
            <pip:description>XNAT User</pip:description>
        </pip:argument>
        <pip:argument id="passwd" isSensitive="true">
            <pip:description>XNAT User password</pip:description>
        </pip:argument>
        <pip:argument id="pipelineLogFile">
            <pip:description>Path to pipeline log file</pip:description>
        </pip:argument>
    </pip:input>
</pip:Resource>
\end{lstlisting}
The resource descriptor XML for an XNAT pipeline specifies the executable's invocation details and argument handling—such as the location, command prefix, and any argument marked as `isSensitive` to hide credentials.\footnote{Mohana Ramaratnam (2013, July 25). “Re: Running pipeline with parameters from an XML?”, *xnat\_discussion* Google Group. Available at: \url{https://groups.google.com/g/xnat_discussion/c/ewhjL7upJf4} (last accessed 10 September 2025).}
The  Listing~\ref{lst:pipeline-xml} presents a sample resource descriptor for a Python script, specifying its location, command prefix, and input arguments such as configuration file, XNAT host, user credentials (with password marked as sensitive), and log file path. This example illustrates how the schema integrates execution details with security considerations for sensitive data.
The commands shown are part of the process of setting up a classic XML-based pipeline in XNAT.
The first command copies the file my\_pipeline.xml, which contains the XML-based definition of the pipeline
workflow, into the directory /data/xnat/pipeline/pipelines/my\_pipeline/, where XNAT expects all pipeline
descriptors to be stored.
The second command copies the file run.sh—a shell script that performs the actual data processing defined in the
pipeline into the same directory. This ensures that both the pipeline definition and the script it executes are
located together in the correct structure required by the XNAT Pipeline Engine. Finally, the third command
makes the run.sh script executable using chmod +x, which is essential because XNAT will later attempt to run
this script as a command. Without executable permissions, XNAT would not be able to start the script, and the pipeline execution would fail.
\footnote{XNAT Documentation (2025): \url{https://wiki.xnat.org/documentation/installing-pipelines-in-xnat} last Accessed 10 September 2025}

\section{Pipeline Development and Integration in XNAT}
On the contrary, building and enabling a pipeline in XNAT involves considerably more manual steps compared to the Docker-based approach.
First we have to enter the path to the pipeline in XNAT under Admin>Pipelines>Register new Pipelines> And then, in the appropriate field, we enter the path. Then we have to build the option run the Pipeline.In order for users to execute the Pipeline from the XNAT interface an administrator must configure the relevant data type to include a pipeline launch action. This involves editing the data type settings under Administer >Data Types and adding a new action as PipelineScreen\_launch\_pipeline—in the Available Report Actions section, with a suitable display label like Build or Run Pipeline.\footnote{\url{https://wiki.xnat.org/documentation/installing-pipelines-in-xnat}last accessed 10.09.2025}

\section{The Pipeline and the REST API}
With the aim of analyzing the pipeline method more deeply, it is worth noting that XNAT provides several REST API endpoints for managing and running pipelines the most relevant endpoint for automation is \url{POST /pipelines/launch/{pipelineName}} \footnote{\url{://wiki.xnat.org/xnat-api/xnat-pipeline-api} last Accessed 10.09.2025}
This allows us to programmatically launch an already registered pipeline in XNAT. However, for full
automation, additional endpoints would be needed, such as one to enable a pipeline globally in XNAT and
another to assign or activate a pipeline for a specific project. These functionalities are currently not exposed via
the REST API, which means that while the execution of pipelines can be automated, the full deployment and
project-level configuration still require manual steps.

\section{Evaluating Pipeline and Docker Approaches in XNAT}

For the purpose of evaluating the two methods, we began by compiling a comparative table to provide a clearer understanding of their respective features. 

\begin{table}[H]
    \centering
      \begin{tabular}{|>{\centering\arraybackslash}p{4cm}|
                    >{\centering\arraybackslash}p{5cm}|
                    >{\centering\arraybackslash}p{5cm}|}
    \hline
    \textbf{Aspect} & \textbf{The Pipeline method} & \textbf{Docker Container method}\\ \hline
     Installation & Requires manual placement of XML descriptors, plugin JARs, and executables on the server& Container image uploaded once (via Docker registry) and configured in XNAT. \\ \hline
    Configuration & Admin must edit data types, enable actions, and rebuild pipeline engine, require direct access to the server via SSH in order to manually copy the XML
descriptor and associated scripts to a specific directory (/data/xnat/pipeline/pipelines/)& Minimal configuration through a JSON command; no engine rebuild needed. \\ \hline
Execution & the Run Pipeline option must be build via Actions menu  & Triggered via Actions menu or REST API, image runs in isolated environment. \\ \hline
Automation & Limited REST API Supported by XNAT Community, many tasks still require manual steps.& Excellent support of the REST API by XNAT Community. Enabling a complete automation of workflows. \\ \hline
Deployability and flexibility & Pipelines are highly dependent on the specific XNAT environment and server configuration in which they are deployed. & Docker images can be readily reused in various XNAT installations.\\ \hline
Preservation & Updating necessitates server access and manual reconfiguration.& Updates can be managed through the construction and deployment of a new image, with no need for direct changes to the server configuration.\\ \hline
    \end{tabular}
    \caption{Comprehensive comparison of XNAT Pipeline Method and Docker Container Service}
    \label{tab:pipeline-vs-docker}
\end{table}
 The table 4.1 has been added to present a direct comparison between the two methods.
\section{Discussion and Conclusion and future measurements}



  %\chapter{User Guide}
To better understand the system created and mostly the motivation behind this automation project. We have inserted a table to list how the the Workflow is working and what are the requirement that the user needs to have in order to run the script.


\begin{center}
\begin{tabular}{|l|p{9cm}|}
\hline
\textbf{System Requirements} & 
- Python installed \newline
- \texttt{requirements.txt} in the same folder as \texttt{automatic.py} \newline
- XNAT server accessible \newline
- Docker Engine installed and running \newline
- Logged in to Docker Hub \\
\hline
\textbf{Script Name} & \texttt{automatic.py} \\
\hline
\textbf{Required Libraries} & requests, json, subprocess, getpass, os, sys, urllib3 \\
\hline
\textbf{User Inputs (Prompted)} & 
- XNAT Username \newline
- XNAT Password \newline
- Project ID \newline
- Docker Hub Username \newline
- Container Script (\texttt{.py}) \newline
- Command Name and Description \newline
- Wrapper Label and Description \\
\hline
\textbf{Steps to Follow} & 
1. Run the \texttt{automatic.py} script \newline
2. Enter the prompted inputs \newline
3. Wait until the container is launched \\
\hline
\textbf{What Happens Automatically} & 
- Dockerfile is generated \newline
- Docker image is built, tagged, and pushed \newline
- Command JSON is generated \newline
- Command is uploaded to XNAT \newline
- Wrapper is activated site-wide and project-wide \newline
- Container is launched via REST API \\
\hline
\textbf{Problems / Troubleshooting} & 
- Push failed: check Docker Hub login \newline
- Wrapper not found or JSON upload failed: check chosen command naming \newline
- Command already exists: delete old command manually \newline
- Container problem uploading files: try running \texttt{ls -al \*} on your output directory and check the container \texttt{stdout} log \\
\hline
\end{tabular}
\end{center}


  %\chapter{Conclusion}


\section{Conclusion}
In this work, a framework was developed and automatized for executing machine learning models in the open-source medical Imaging and biosignals platform XNAT. This Paper presented the development and the Debugging process involved in integrating a Docker Container within the XNAT Environment for automated data processing.

Numerous Challenges made an appearance, first and foremost related to file input handling, container execution context, and the behavior of XNAT's data mounting system. Throughout repetitive testing and including diagnostic Scripting, REST API Integration, and JSON Command modifications, in spite of the fact that the core failures became comprehensible, a resolution to the underlying issue has not been achieved.


In Reply to These challenges, several countermeasures were proposed. These comprise implementing alternative procedures in Container scripts, ameliorating the Container logs for a better diagnostics. Furthermore, the analysis exposed deficiencies in XNAT’s temporary data handling and its limited error reporting capabilities, emphasizing the need for targeted improvements within the platform.

In a wider context, this work underline the complexity of the practice of integrating containers into clinical Research platforms.
Despite encountering technical constraints, the findings gained throughout the process
contribute valuable knowledge for future container based workflows in XNAT. The scripts and procedures developed form a pivotal step in establishing data processing pipelines that offer increased reliability and scalability. This work contributes meaningfully to the field of medical informatics, through this approach, it can be used in automated patient data processing, medical AI Workflow. It also provides a fast and efficient way to create new operations or updates, which can have positively a huge impact of the patient diagnosis or on the patient care in general.

The next phase involves a more in-depth examination of XNAT's mechanisms for handling file mounting at the system Level. Collaborating with XNAT developers or administrators may also help identify undocumented behaviors or system constraints. Additionally, exploring alternative methods for data Transfer such as Volume mounting from external storage or utilizing wrapper containers supporting the management of multiple files  might strengthen overall the results. Finally, transitioning from per-file execution to grouped workflows and container orchestration strategies could markedly increase scalability while minimizing processing demands in extensive project environments.




  %--- references ---
  \bibliographystyle{IEEEtran}
  \bibliography{content/references}
  \addcontentsline{toc}{chapter}{\bibname}

  %--- appendix ---
  \begin{appendix}
    \appendix
\section{Test Protocol and Results}
\label{app:test}

\noindent\textbf{Author:} Tanae Bousfiha  
\noindent\textbf{Date:} August 2025  

\noindent This appendix documents the full terminal output of the automated Docker–XNAT integration script
\texttt{Projekt.py}.  
It includes the complete build, push, and wrapper activation process, as well as the execution log showing the files processed by the container.  
This section serves as a reproducibility record for the workflow described in Section~X.X of the main text.

\begin{lstlisting}
tanae@mi0082n:~/dockerfile/xnat-osa/Ziel2_osaAutomatisation$ python Projekt.py
XNAT Username: tanae
XNAT Password:
Project ID:  tanae_osa_Predct
Path to the Python script: test.py
Name des Commands: Realisation
Beschreibung des Commands: describes the Realistaion Process
Dockerfile written to ./Dockerfile
Docker Hub username (to push the image): tanaebousfiha
Building Docker image 'realistaion:latest'...
Image 'realistaion:latest' built successfully.
Tagging image as 'tanaebousfiha/realistaion:latest'...
Pushing image to Docker Hub as 'tanaebousfiha/realistaion:latest'...
Image successfully pushed: tanaebousfiha/realistaion:latest
Corrected command.json created.
Uploading command to https://xnat-dev.gwdg.de/xapi/commands
Command created successfully.
Command-ID: 637
Wrapper-ID: 674
Wrapper 'realistaion_wrapper' wurde global aktiviert.
Wrapper 'realistaion_wrapper' wurde im Projekt 'tanae_osa_Predct' aktiviert.
Start the Container.:
Dateien: ['test.csv', 'DemandPlan_v1.csv', 'Inventory_v2.csv', ... more than 700 files]
Container succesfully launched with all files!
\end{lstlisting}

  \end{appendix}

\end{document}

